\section*{Chapter \ref{ch:m3:natural-gas-and-lng} --- Natural Gas and LNG}

\begin{solution}{exo:m3:natural-gas-and-lng:1}
$10/0.293071/1.16 = \qty{29.42}{EUR/MWh}$. A therm is 0.1 MMBtu, so \$1 a therm, or
$100/1.30 = 76.9$ pence per therm.
\end{solution}

\begin{solution}{exo:m3:natural-gas-and-lng:2}
Winter demand exceeds production and must be met from storage filled in summer, so
winter gas must pay for storing summer gas. Arbitrage by storage owners bounds it from
above by the cost of storage (injection, withdrawal, capacity and financing) when
capacity is free; when storage is full or scarce it can exceed that cost.
\end{solution}

\begin{solution}{exo:m3:natural-gas-and-lng:3}
The variable price is $1.15 \times 2.00 = \qty{2.30}{\$/\mmbtu}$, above the \omterm{def:m3:physical-commodity-markets:netback}{netback} of
2.10: it cancels. It loses the \qty{2.50}{} fee either way; lifting would lose a further
\qty{0.20}{\$/\mmbtu}.
\end{solution}

\begin{solution}{exo:m3:natural-gas-and-lng:4}
Above $11.00 + 0.33 = \qty{11.33}{\$/\mmbtu}$.
\end{solution}

\begin{solution}{exo:m3:natural-gas-and-lng:5}
The holders of pipeline capacity out of the hub: the spread above transport cost,
\qty{1.70}{\$/\mmbtu}, is the value of their capacity. With capacity, buy at the hub, ship,
and sell at \omterm{def:m3:natural-gas-and-lng:hubs}{Henry Hub} (or sell the basis forward).
\end{solution}

\begin{solution}{exo:m3:natural-gas-and-lng:6}
The break-even spread rises from \qty{0.33}{} to \qty{1.35}{\$/\mmbtu}, by \qty{1.02}{}: the
longer voyage costs more charter and more boil-off.
\end{solution}

\begin{solution}{exo:m3:natural-gas-and-lng:7}
Asia in 93 of the 125 months from March 2016 to July 2026; not lifted in February to August
2020 and February 2021.
\end{solution}

\begin{solution}{exo:m3:natural-gas-and-lng:8}
The cargo's revenue is the destination price less shipping and regasification, and its
cost is 115\% of \omterm{def:m3:natural-gas-and-lng:hubs}{Henry Hub} plus the fixed fee paid for the capacity: the margin is the
\omterm{def:m3:physical-commodity-markets:netback}{netback} less both. Most of the spread was also earned by whoever held the cargo's
destination rights, often under long-term contracts signed years earlier, not by every
seller; and freight rates rose with the spread.
\end{solution}

\begin{solution}{pb:m3:natural-gas-and-lng:1}
\textbf{1.} \qty{3.45}{\$/\mmbtu}. \textbf{2.} \qty{0.54}{\$/\mmbtu}. \textbf{3.}
\qty{0.97}{\$/\mmbtu}. \textbf{4.} 1.4\% to Europe, 2.5\% to Asia. \textbf{5.} The ship
must come back empty (in ballast) before it can load again.
\textbf{6.} \qty{0.33}{\$/\mmbtu}. \textbf{7.} Asia: \omterm{def:m3:physical-commodity-markets:netback}{netback} \qty{11.13}{} against
\qty{10.96}{}. \textbf{8.} $11.13 \times 3.4125$ million delivered MMBtu less $3.45 \times
3.5$ million: about USD~25.9 million. \textbf{9.} The break-even rises to \qty{0.72}{}, above
the \qty{0.50}{} spread: Europe. \textbf{10.} Europe: \omterm{def:m3:physical-commodity-markets:netback}{netback} \qty{68.85}{} against
\qty{52.62}{\$/\mmbtu}.
\textbf{11.} Prices at the two destinations move apart at random; the right to pick the
higher, net of the break-even, is an option on the spread, worth something whenever the
spread is volatile. \textbf{12.} Not to lift: the best \omterm{def:m3:physical-commodity-markets:netback}{netback}, \qty{0.72}{\$/\mmbtu}, was
below the variable price of \qty{1.87}{}. \textbf{13.} Nothing: the fee is owed whether or
not the cargo is lifted, so it is not part of the marginal decision. \textbf{14.} The
buyer (it saves the variable price minus the \omterm{def:m3:physical-commodity-markets:netback}{netback}); the liquefaction owner still
receives the fee. \textbf{15.} Freight: charter rates rise, which narrows the arbitrage.
\textbf{16.} Its hubs paid more than Asia after the loss of Russian pipeline gas, and the
voyage from the US Gulf is shorter. \textbf{17.} Asian cargoes from the Atlantic must sail
the longer route, which raises the break-even and so the premium Asia must pay.
\textbf{18.} For price stability tied to its other costs, security of supply, or because
oil-linked contracts were what the seller offered. \textbf{19.} \emph{Named result:} Asia
must pay \qty{0.33}{\$/\mmbtu} more than Europe (\qty{1.35}{} if the voyage takes 50 days).
\textbf{20.} Flexible LNG cargoes, which move to the best \omterm{def:m3:physical-commodity-markets:netback}{netback} until the spreads fall to
the cost of shipping.
\end{solution}

\begin{solution}{iq:m3:natural-gas-and-lng:1}
\omterm{def:m3:natural-gas-and-lng:hubs}{Henry Hub} is a physical interconnection of pipelines in Louisiana and the delivery point
of the US futures; TTF is a \omterm{def:m3:natural-gas-and-lng:hubs}{virtual trading point} of the Dutch network where any gas in
the system can change hands. Different units and currencies too.
\iqlookfor{physical versus virtual hub, and the units.}
\end{solution}

\begin{solution}{iq:m3:natural-gas-and-lng:2}
Gas demand is dominated by heating and power, both seasonal, while production is flat
and storage costly; oil demand is less seasonal and oil is cheap to store in tanks.
\iqlookfor{seasonal demand against flat supply and storage cost.}
\end{solution}

\begin{solution}{iq:m3:natural-gas-and-lng:3}
Each month, an option to buy gas at 115\% of \omterm{def:m3:natural-gas-and-lng:hubs}{Henry Hub} and sell it at the best
destination hub less shipping and regasification: a spread option (on the maximum of
several spreads), with the fixed fee as its premium.
\iqlookfor{spread option, variable price as strike, fee as premium.}
\end{solution}

\begin{solution}{iq:m3:natural-gas-and-lng:4}
A basis blow-out: a pipeline outage or a local demand spike moves the local hub far from
\omterm{def:m3:natural-gas-and-lng:hubs}{Henry Hub}, beyond transport cost; weather events (a freeze) move both legs by different
amounts; margin calls on the paper legs.
\iqlookfor{basis and capacity, not the price level.}
\end{solution}

\begin{solution}{iq:m3:natural-gas-and-lng:5}
They had sold their output forward on exchanges; when prices rose tenfold the short
futures required variation margin in cash, while the physical gas that offset them would
only be sold later. The hedge was economically sound but consumed liquidity they did not
have.
\iqlookfor{liquidity versus solvency, margin on exchange hedges.}
\end{solution}

\begin{solution}{iq:m3:natural-gas-and-lng:6}
Inputs: cargo positions and schedules, contractual constraints, forward curves per hub,
freight curves and canal status. For each cargo, compute the \omterm{def:m3:physical-commodity-markets:netback}{netback} to each feasible
destination and date, choose subject to terminal slots and ship availability (an
assignment problem), and report the choice, the margin and its sensitivity to spreads.
\iqlookfor{\omterm{def:m3:physical-commodity-markets:netback}{netbacks}, constraints, an assignment optimisation.}
\end{solution}
